\documentclass[11pt]{article}
\usepackage[T1]{fontenc}
\usepackage{lmodern}
\usepackage[a4paper,margin=27mm]{geometry}
\usepackage{amsmath,amssymb,amsthm}
\usepackage{graphicx}
\usepackage{microtype}
\usepackage[numbers,sort&compress]{natbib}
\usepackage{hyperref}
\hypersetup{colorlinks=true,linkcolor=blue,citecolor=blue,urlcolor=blue,
  pdftitle={Time-independent Lindblad embeddings of two-parameter Werner-Holevo noise},
  pdfauthor={Xiaojian Zeng}}
\newcommand{\id}{\operatorname{id}}
\newcommand{\tr}{\operatorname{Tr}}
\newcommand{\Ad}{\operatorname{Ad}}
\newcommand{\OO}{\mathrm{O}}
\newcommand{\SO}{\mathrm{SO}}
\newcommand{\Ps}{\mathsf{P}_{\mathrm{s}}}
\newcommand{\Pa}{\mathsf{P}_{\mathrm{a}}}
\newcommand{\Pz}{\mathsf{P}_{0}}
\newcommand{\gs}{\gamma_{\mathrm{s}}}
\newcommand{\ga}{\gamma_{\mathrm{a}}}
\newtheorem{theorem}{Theorem}
\newtheorem{corollary}[theorem]{Corollary}
\newtheorem{lemma}[theorem]{Lemma}
\renewcommand{\topfraction}{0.9}
\renewcommand{\textfraction}{0.05}
\renewcommand{\floatpagefraction}{0.85}
\title{Time-independent Lindblad embeddings of\\
two-parameter Werner--Holevo noise}
\author{Xiaojian Zeng\\[2pt]
\small PBC School of Finance, Tsinghua University, Beijing 100190, China\\
\small \texttt{xiaojianzeng@foxmail.com}}
\date{}
\begin{document}
\maketitle

\begin{abstract}
An orthogonally covariant quantum channel need not have an orthogonally
covariant time-independent generator. This distinction matters when deciding
whether a calibrated noise channel can arise from one constant Lindblad
master equation. We consider the two-parameter mixture of the identity and
the symmetric and antisymmetric Werner--Holevo channels on a finite
$d$-level system. For positive eigenvalues $s$ and $a$ on the symmetric
traceless and antisymmetric operator sectors, respectively, we prove that
an arbitrary same-system Lindblad embedding exists exactly when
$s^{(d+2)/d}\leq a\leq s^{(d-2)/d}$.
The covariant generator and its rate cone are established prior results;
the additional step is necessity without imposing covariance, a
Hamiltonian restriction, or a logarithm branch. Real sector traces fix the
average logarithmic rates, and averaging over the full orthogonal group
then produces the principal covariant generator. Odd sector multiplicities
exclude every remaining signed qutrit case, yielding a complete
classification for $d=3$. We give normalized jump operators and illustrate
the qutrit region. Negative even-multiplicity sectors in higher dimensions
remain outside the classification.
\end{abstract}
\noindent\textbf{Keywords:} quantum channels; Lindblad embedding;
Werner--Holevo noise; orthogonal covariance; quantum Markov semigroups

\section{Endpoint noise and a constant master equation}

A quantum channel at one finite time does not by itself specify a
time-independent master equation. The endpoint embedding problem asks whether
$\Phi=e^{T\mathcal L}$ for a Gorini--Kossakowski--Lindblad--Sudarshan (GKLS)
generator $\mathcal L$ on the same system. General logarithm and conditional
complete-positivity tests are available \cite{Wolf2008}, but degeneracies
allow substantial freedom in the logarithm. In particular, symmetry of an
endpoint cannot simply be imposed on an unknown generator.

The noisy Werner--Holevo models provide a concrete setting for this issue
\cite{Karimipour2024,Roofeh2025}. Their decomposition into orthogonal operator
sectors is established, including in real randomized benchmarking
\cite{Hashagen2018}. Bhat, Chakraborty and Franz studied the generator family
spanned by depolarization, transposition and the identity, and derived its
semigroup exponential and positivity conditions \cite{Bhat2023}. Thus the
covariant family and the associated rate cone are already known. What remains
to justify for an endpoint test is that an arbitrary, possibly noncovariant
GKLS logarithm cannot evade that cone.

Here we prove this reduction for the positive-spectrum part of the
two-parameter model. The argument uses real traces on degenerate endpoint
sectors rather than an enumeration of generic-spectrum logarithm branches.
It gives a complete endpoint classification for qutrits after the signed
sectors are excluded. Our result concerns a single finite, time-independent
generator; it does not classify time-dependent completely positive divisible
paths. The related Pauli-channel semigroup criteria of
Kattem\"olle, Gul\'acsi and Burkard
(\href{https://arxiv.org/abs/2602.08464v1}{arXiv:2602.08464v1}) concern a different
operator algebra and supply a useful conceptual comparison.

\section{Model, sectors and physical domain}

Fix an orthonormal computational basis of $\mathbb C^d$, with $d\geq3$.
Transposition below is the complex-linear map $X\mapsto X^{\mathsf T}$ in
this basis, not Hermitian conjugation. Following
Ref.~\cite{Roofeh2025}, define
\begin{align}
 W_-(X)&=\frac{\tr(X)I_d-X^{\mathsf T}}{d-1},&
 W_+(X)&=\frac{\tr(X)I_d+X^{\mathsf T}}{d+1},\label{eq:WH}\\
 \Phi_{d,x,y}&=(1-x-y)\id+xW_-+yW_+,\label{eq:model}\\
 &x>0,\quad y>0,\quad x+y\leq1. \nonumber
\end{align}
All inputs in $M_d(\mathbb C)$ are included. In particular, complete
positivity is required on entangled inputs. No enlarged system, postselection
or approximation is allowed in the embedding.

Set
\begin{equation}
 q=1-x-y,\qquad c=\frac{x}{d-1}+\frac{y}{d+1},\qquad
 b=\frac{y}{d+1}-\frac{x}{d-1}.
\end{equation}
Then $\Phi(X)=qX+c\,\tr(X)I_d+bX^{\mathsf T}$. The orthogonal
superoperator projectors are
\begin{equation}
 \Pz(X)=\frac{\tr(X)}dI_d,\qquad
 \Ps(X)=\frac{X+X^{\mathsf T}}2-\Pz(X),\qquad
 \Pa(X)=\frac{X-X^{\mathsf T}}2.                 \label{eq:projectors}
\end{equation}
Their complex dimensions, also their dimensions on real Hermitian bases, are
\begin{equation}
 1,\qquad m_{\mathrm{s}}=\frac{d(d+1)}2-1,\qquad
 m_{\mathrm{a}}=\frac{d(d-1)}2.
\end{equation}
The traceless Hermitian bases may be chosen as real symmetric matrices and
purely imaginary antisymmetric matrices. The channel is
\begin{align}
 \Phi&=\Pz+s\Ps+a\Pa,                                      \label{eq:sectors}\\
 s&=q+b=1-\frac{d}{d-1}x-\frac{d}{d+1}y, \nonumber\\
 a&=q-b=1-\frac{d-2}{d-1}x-\frac{d+2}{d+1}y. \label{eq:eigenvalues}
\end{align}
Both $s$ and $a$ are strictly less than one in the domain
\eqref{eq:model}; their signs are not fixed by complete positivity.

For clarity, use the unnormalized Choi convention
\begin{equation}
 J(\mathcal M)=\sum_{i,j=1}^d E_{ij}\otimes\mathcal M(E_{ij}),\qquad
 |\Omega\rangle=\sum_{i=1}^d|ii\rangle,
\end{equation}
with input followed by output. If $F$ is the swap, then
$J(\Phi)=q|\Omega\rangle\langle\Omega|+cI+bF$.
Its eigenvalues are
\begin{equation}
 dq+\frac{2y}{d+1},\qquad
 \frac{2y}{d+1},\qquad \frac{2x}{d-1},             \label{eq:choichannel}
\end{equation}
with multiplicities $1,m_{\mathrm{s}},m_{\mathrm{a}}$.
They are strictly positive even on $x+y=1$, and
$\tr_{\mathrm{out}}J(\Phi)=I_d$. A full-rank Choi matrix therefore does
not imply an invertible channel superoperator or a Lindblad embedding.

We set $T=1$ and $\hbar=1$. A generator has the form
\begin{equation}
 \mathcal L(X)=-i[H,X]
 +\sum_r\left(V_rXV_r^\dagger-\frac12\{V_r^\dagger V_r,X\}\right),
 \qquad H=H^\dagger.                              \label{eq:GKLS}
\end{equation}
For a general observation time $T>0$, divide the generator and its jump
rates below by $T$, and divide the displayed jumps by $\sqrt T$.
The endpoint weights and eigenvalues are dimensionless.

\section{Unrestricted positive-spectrum criterion}

\begin{theorem}\label{thm:positive}
For every integer $d\geq3$, let $\Phi_{d,x,y}$ satisfy \eqref{eq:model}
and suppose $s>0$ and $a>0$. There exists a finite GKLS generator
$\mathcal L$ on $M_d(\mathbb C)$ with $e^{\mathcal L}=\Phi_{d,x,y}$
if and only if
\begin{equation}
             s^{(d+2)/d}\leq a\leq s^{(d-2)/d}.     \label{eq:cone}
\end{equation}
Necessity allows arbitrary Hamiltonians, jump operators, logarithm branches
and transformations within degenerate endpoint eigenspaces.
A generator can be chosen with $H=0$ and
$\mathcal L_0=(\log s)\Ps+(\log a)\Pa$.
\end{theorem}

The cone in \eqref{eq:cone} is the endpoint expression of the known
covariant-generator cone \cite{Bhat2023}. The content of the theorem is
that the same cone is necessary for every generator of this endpoint.

\subsection{Real sector traces and the orthogonal average}

First assume $s\ne a$. Since any matrix commutes with its exponential,
$\mathcal L$ commutes with $\Phi$. The distinct eigenvalues $1,s,a$
force $\mathcal L$ to preserve the three sectors in \eqref{eq:projectors}.
On the scalar sector $\mathcal L(I_d)=hI_d$; trace preservation gives
$h=0$. Hermiticity preservation makes the restrictions
$\mathcal L_{\mathrm{s}}$ and $\mathcal L_{\mathrm{a}}$ real matrices
in their Hermitian bases. Consequently their traces are real and
\begin{align}
 e^{\tr\mathcal L_{\mathrm{s}}}
 &=\det e^{\mathcal L_{\mathrm{s}}}=s^{m_{\mathrm{s}}},&
 \tr\mathcal L_{\mathrm{s}}&=m_{\mathrm{s}}\log s,\nonumber\\
 e^{\tr\mathcal L_{\mathrm{a}}}
 &=\det e^{\mathcal L_{\mathrm{a}}}=a^{m_{\mathrm{a}}},&
 \tr\mathcal L_{\mathrm{a}}&=m_{\mathrm{a}}\log a.  \label{eq:traces}
\end{align}
These identities do not select a logarithm branch or assume normality of
the restrictions.

The endpoint has full orthogonal covariance:
$\Ad_O(X)=OXO^{\mathsf T}$ commutes with $\Phi$ for every $O\in\OO(d)$.
Average the generator, with normalized Haar measure,
\begin{equation}
 \overline{\mathcal L}=
 \int_{\OO(d)}\Ad_O^{-1}\mathcal L\Ad_O\,dO.        \label{eq:average}
\end{equation}
Unitary conjugation preserves GKLS form. The finite-dimensional GKLS cone
is closed and convex: trace and Hermiticity preservation are linear
conditions and conditional complete positivity is a linear matrix
inequality \cite{Wolf2008,Bhat2023}. Thus $\overline{\mathcal L}$ is
again GKLS. It preserves the sectors, vanishes on the scalar sector, and
retains their restricted traces.

We spell out why its restrictions are scalar for the full group. Diagonal
sign matrices act on each off-diagonal pair $(i,j)$ with character
$\varepsilon_i\varepsilon_j$. These characters are distinct for the full
set of independent signs, so an intertwiner within either sector scales
each off-diagonal basis vector. Permutation matrices make those scales
equal. In the symmetric sector, the diagonal traceless subspace is the
sum-zero permutation representation; an intertwiner is scalar there as
well. A planar rotation by $\pi/4$ maps $E_{ii}-E_{jj}$ to a signed
multiple of $E_{ij}+E_{ji}$, equating the diagonal and off-diagonal
symmetric scales. The antisymmetric sector has no diagonal component.
This elementary argument gives a scalar on each sector for all $d\geq3$.
In $d=4$, using full $\OO(4)$ is essential: reflections exchange the two
Hodge components that are separately invariant under $\SO(4)$.
The corresponding orthogonal-sector decomposition is also described in
Ref.~\cite{Hashagen2018}.

Equation \eqref{eq:traces} now fixes the two scalar restrictions:
\begin{equation}
       \overline{\mathcal L}=(\log s)\Ps+(\log a)\Pa
       =\mathcal L_0.                              \label{eq:reduction}
\end{equation}
Therefore an arbitrary GKLS embedding implies that the principal sector
logarithm is GKLS. This conclusion does not follow from a general
commutation of averaging and exponentiation; \eqref{eq:reduction} follows
from the specific real trace identities.

When $s=a>0$, the distinct-sector argument is unnecessary. The channel is
depolarizing and
$\mathcal L_0=(\log s)(\id-\Pz)$ is an explicit GKLS generator.
Because $0<s<1$, \eqref{eq:cone} holds automatically. This handles the
coincident traceless eigenvalues without assuming sector preservation for
an arbitrary logarithm on that line.

\subsection{The known rate cone in endpoint variables}

Let $\ell_{\mathrm{s}}=\log s$, $\ell_{\mathrm{a}}=\log a$ and
$Q=I-|\Omega\rangle\langle\Omega|/d$. Direct Choi reshuffling yields
\begin{align}
 J(\mathcal L_0)
 &=\frac{\ell_{\mathrm{s}}+\ell_{\mathrm{a}}}{2}
       |\Omega\rangle\langle\Omega|
   +\frac{\ell_{\mathrm{s}}-\ell_{\mathrm{a}}}{2}F
   -\frac{\ell_{\mathrm{s}}}{d}I, \nonumber\\
 QJ(\mathcal L_0)Q&=\gs C_{\mathrm{s}}+\ga C_{\mathrm{a}},       \label{eq:ccp}\\
 C_{\mathrm{s}}&=\frac{I+F}{2}
                   -\frac{|\Omega\rangle\langle\Omega|}{d},
 &C_{\mathrm{a}}&=\frac{I-F}{2},\nonumber\\
 \gs&=\frac{(d-2)\ell_{\mathrm{s}}-d\ell_{\mathrm{a}}}{2d},
 &\ga&=\frac{d\ell_{\mathrm{a}}-(d+2)\ell_{\mathrm{s}}}{2d}.
                                                               \label{eq:rates}
\end{align}
The Choi-sector projectors $C_{\mathrm{s}},C_{\mathrm{a}}$ are mutually
orthogonal and distinct from the superoperator notation
$\Ps,\Pa$. Conditional complete positivity is exactly
$\gs,\ga\geq0$, which is equivalent to \eqref{eq:cone}.
This proves necessity for $s\ne a$.

To make the prior-result correspondence explicit, write
$\mathcal T(X)=X^{\mathsf T}$. Section 6.3 of Ref.~\cite{Bhat2023} uses
\begin{equation}
 \mathcal G_{\alpha,\eta}
  =\alpha\Pz+\eta\mathcal T-(\alpha+\eta)\id.
 \label{eq:prior}
\end{equation}
The exact parameter map is
\begin{equation}
 \alpha=-\ell_{\mathrm{s}},\qquad
 \eta=\frac{\ell_{\mathrm{s}}-\ell_{\mathrm{a}}}{2},\qquad
 \gs=\frac{\alpha}{d}+\eta,\qquad
 \ga=\frac{\alpha}{d}-\eta.                        \label{eq:map}
\end{equation}
Thus its covariant generator cone $\alpha\geq d|\eta|$ is precisely
\eqref{eq:cone}. The symbol $\eta$ denotes the parameter called $\gamma$
in that reference, to avoid confusing it with the two jump rates here.
Its exponential has sector eigenvalues
$e^{-\alpha t}$ and $e^{-(\alpha+2\eta)t}$.
Neither this exponential nor its covariant rate cone is claimed as a new
result.

\subsection{Normalized jumps and sufficiency}

Choose Frobenius-orthonormal Hermitian bases
$\{S_\mu\}_{\mu=1}^{m_{\mathrm{s}}}$ of real symmetric traceless matrices
and $\{A_\nu\}_{\nu=1}^{m_{\mathrm{a}}}$ of purely imaginary
antisymmetric matrices. Their completeness identities are
\begin{align}
 \sum_\mu S_\mu XS_\mu^\dagger
 &=\frac{\tr(X)I_d+X^{\mathsf T}}2-\frac Xd,
 &\sum_\mu S_\mu^\dagger S_\mu&=\frac{m_{\mathrm{s}}}{d}I_d,\nonumber\\
 \sum_\nu A_\nu XA_\nu^\dagger
 &=\frac{\tr(X)I_d-X^{\mathsf T}}2,
 &\sum_\nu A_\nu^\dagger A_\nu&=\frac{m_{\mathrm{a}}}{d}I_d.
                                                            \label{eq:completeness}
\end{align}
For example, off-diagonal basis elements are
$(E_{ij}+E_{ji})/\sqrt2$ and $i(E_{ij}-E_{ji})/\sqrt2$.
A diagonal symmetric basis is
\begin{equation}
 S^{\mathrm{diag}}_k=
 \frac{\sum_{i=1}^k E_{ii}-kE_{k+1,k+1}}{\sqrt{k(k+1)}},
 \qquad 1\leq k\leq d-1.
\end{equation}
If \eqref{eq:cone} holds, take $H=0$ and jumps
\begin{equation}
       V_\mu=\sqrt{\gs}\,S_\mu,\qquad
       V_\nu=\sqrt{\ga}\,A_\nu.                   \label{eq:jumps}
\end{equation}
Substituting \eqref{eq:completeness} into \eqref{eq:GKLS} gives exactly
$\mathcal L_0=\ell_{\mathrm{s}}\Ps+\ell_{\mathrm{a}}\Pa$.
Orthogonality of the superoperator projectors gives
\begin{equation}
 e^{t\mathcal L_0}
  =\Pz+s^t\Ps+a^t\Pa, \qquad
 e^{\mathcal L_0}=\Phi.                           \label{eq:exp}
\end{equation}
This proves sufficiency and completes Theorem~\ref{thm:positive}.
On the two boundaries of the cone, one of the jump families has zero
rate; a finite generator still exists.

\section{Complete qutrit classification}

\begin{corollary}\label{cor:qutrit}
For $d=3$ throughout $x>0,y>0,x+y\leq1$, a finite time-independent
same-qutrit GKLS embedding exists exactly when
\begin{equation}
 s=1-\frac32x-\frac34y>0,\qquad
 a=1-\frac12x-\frac54y>0,\qquad
                    s^5\leq a^3\leq s.           \label{eq:qutrit}
\end{equation}
\end{corollary}
\begin{proof}
The two real Hermitian sector multiplicities are $m_{\mathrm{s}}=5$
and $m_{\mathrm{a}}=3$. A zero eigenvalue makes $\Phi$ singular, which
is impossible for a finite matrix exponential. If one eigenvalue is
negative, then $s+a=2q\geq0$ prevents both from being negative; the
nonzero eigenvalues are therefore distinct and the sector restrictions
of an embedding would be real. A real matrix exponential has positive
determinant. A negative scalar on either odd-dimensional sector has
negative determinant, so no such logarithm exists. All remaining
possibilities have $s,a>0$ and Theorem~\ref{thm:positive} gives
\eqref{eq:qutrit}.
\end{proof}

\begin{figure}[t]
\centering
\includegraphics[width=0.82\textwidth]{manuscript/figures/qutrit-region.pdf}
\caption{Exact homogeneous-embedding region in the qutrit noise triangle.
The colored region includes its two rate boundaries wherever $x,y>0$;
the axes and the limiting point $(1/3,2/3)$ are excluded.
The boundary $\gs=0$ has $a^3=s$, and $\ga=0$ has $a^3=s^5$.
The dotted line $y=2x$ is the depolarizing subfamily.
The circle at $(1/4,1/4)$ is embeddable, whereas the cross at
$(1/2,1/10)$ has positive spectrum but violates the cone.
Curves are sampled from the analytic boundary equations for rendering;
the classification follows from Corollary~\ref{cor:qutrit}.}
\label{fig:region}
\end{figure}

Figure~\ref{fig:region} separates full Choi rank from homogeneous
embeddability. At $x=y=1/4$, the endpoint parameters are
$s=7/16$ and $a=9/16$, and both rate inequalities are strict.
This same input reproduces the published output benchmark in
Ref.~\cite{Roofeh2025}: the pure state
$|\psi\rangle=(i,1,0)^{\mathsf T}/\sqrt2$ has output spectrum
$\{1/8,3/16,11/16\}$. This is a normalization check, rather than a new
capacity result.

In contrast, the legal point $x=1/2,y=1/10$ gives
$s=7/40$ and $a=5/8$, with
$a^3=125/512>s$. Its principal symmetric jump rate is negative.
Theorem~\ref{thm:positive} shows that a Hamiltonian, noncovariant jumps
or a different logarithm branch cannot repair that failure.
The statement concerns the exact channel at a specified time; applying it
to a noisy experimental estimate would require uncertainty propagation
and an independent model-validity assessment.

\section{Interpretation and limitations}

The result converts a known covariant-generator cone into an unrestricted
endpoint test for the positive orthogonal sectors. This distinction is
physically useful when the measured or prescribed endpoint has symmetry but
the underlying master equation is not assumed to share it. The construction
uses a zero Hamiltonian, while the exclusion statement allows any
Hamiltonian and any finite set of jumps on the same $d$-level system.
It is exact in finite dimension, with no perturbative, continuum or
long-time approximation.

The arbitrary-generator reduction uses positivity of $s$ and $a$.
For an even-dimensional negative sector, a real logarithm may have trace
$m\log|\lambda|$. An orthogonal average would then exponentiate to the
absolute-spectrum endpoint, not necessarily to the original signed
channel. Thus the present proof does not decide negative
even-multiplicity sectors for $d\geq4$. The full signed classification
in arbitrary dimension remains unresolved. The criterion applies to a finite
generator on the original system and does not characterize time-dependent
divisible paths.

The orthogonal sectors are established in Ref.~\cite{Hashagen2018}, the
covariant cone and exponential in Ref.~\cite{Bhat2023}, and the general
conditional-positivity test in Ref.~\cite{Wolf2008}. The additional argument
is the branch-independent real-trace reduction in
\eqref{eq:traces}--\eqref{eq:reduction} and its qutrit consequence.
Priority is limited to the compared literature; no claim of a comprehensive
classification of symmetric quantum-channel embeddings is made.

\clearpage
\section*{Data, code and generative-AI disclosure}

The proof is analytic. The accompanying source package contains exact
matrix-unit and Kraus reconstructions, symbolic Choi and jump checks,
and code and rational coordinate data for Fig.~\ref{fig:region}.
Its reproducibility instructions specify the environments and distinguish
exact identities from optional floating-point diagnostics. No experimental
dataset is used.

OpenAI Codex was used in literature searching and comparison, formulation
of the reduction, derivations, symbolic and numerical code, checking,
and manuscript drafting. These uses extend beyond copy editing.
Separate AI review passes checked the scientific argument, but do not
constitute external peer review. Final scientific review, authorship
confirmation and accountability by the named human author remain pending
for this local draft.

\bibliographystyle{unsrtnat}
\bibliography{manuscript/references}
\end{document}
